The map in the answer sheet represents the sky in azimuthal equidistant projection. The Zenith is at the center of the chart.
\ioaafig{2024_OBS-SKY_OC01-f1.pdf}
\begin{enumerate}
\item Estimate the latitude of the location from where this sky map is captured. \hfill[4 points]
\item Label the cardinal points (N, E, S, W) on the map and mark the three planets from the Solar System in the sky map with a small square around each of them. \hfill[7 points]
\item Draw on the map with continuous lines the celestial equator (EQ), the ecliptic (EC), and the local meridian (ME) and label them. \hfill[4.5 points]
\item Trace 10 constellations in the sky map and label them as C1, C2, \ldots, C10 and write the corresponding Latin name of the constellations or IAU code. \hfill[10 points]
\item Identify and circle the 4 brightest stars in the given sky map. Label the star from S1, the brightest, and continue with the fainter ones till number S4 for the faintest. Write the Bayer designation or the official IAU name of the four identified stars. \hfill[6 points]
\item Mark (with a small circle) positions of the following objects and label them. \hfill[10.5 points]
\[\mathrm{M42},\quad\mathrm{M81},\quad\mathrm{M41},\quad\mathrm{M31},\quad\gamma\ \mathrm{Ori},\quad\alpha\ \mathrm{Per},\quad\alpha\ \mathrm{Cyg}.\]
\item Estimate the equatorial coordinates of the star Aldebaran ($\alpha$ Tau). \hfill[3 points]
\end{enumerate}
% source variant: The standalone observation paper instead gives 43.5 points for OC1 and describes a sky somewhere in Nepal. This transcription follows the 45-point combined paper that matches the official answer sheets and solution.
